% SPDX-License-Identifier: CC-BY-4.0
% Copyright (c) 2026 Martin Schröder <info@swedishembedded.com>

\section{Reproducibility}\label{sec:repro}

Every table and macro is generated by \code{scripts/build\_data.py} from \code{data/runs.csv} (one row per measured run: command, counts, hardware, backend, and the
revisions of the three repositories) and \code{data/claims.csv} (the status of each claim and the run that supports it), and from the run reports in \code{data/runs/}.
The script refuses a claim marked met without a run that exists, a run whose report disagrees with its row, and a pending claim that names evidence. \code{make} builds the PDF.

The runs are made with the \code{voice-loop} sample of the repository, for example
\begin{small}
\begin{verbatim}
voice-loop --backend cuda turns --stream --recordings c0-questions \
    --out-dir answers --base Qwen3-8B --report report.json
\end{verbatim}
\end{small}
using models fetched with \code{brain pull} (Qwen3-8B, Qwen3-TTS-12Hz-0.6B-Base, nemotron-3.5-asr-streaming-0.6b). The sentence and question sets are in the
repository's speech resources with their digests. Revisions, hardware and the state of the card are recorded in each run's row.

The teaching studies are run with
\begin{small}
\begin{verbatim}
voice-loop --backend cuda teach-eval --base Qwen3-8B \
    --scenarios data/teaching/scenarios-test.jsonl --report teaching-test.json
voice-loop --backend cuda teach-read --base Qwen3-8B \
    --utterances data/teaching/utterances.jsonl --report teaching-read.json
\end{verbatim}
\end{small}
on sets written by \code{scripts/teaching\_set.py} and \code{scripts/teaching\_utterances.py}, whose digests are stored beside them; the pre-registrations are in \code{prereg/}.
